\documentclass[12pt]{article}
\usepackage{amsmath, amssymb, amsthm}
\usepackage{graphicx}
\title{Discrete Math HW7}
\author{Anatolii Larkin}

\begin{document}

\maketitle


\begin{enumerate}
\item Solve the system of linear equations using the Gaussian method.
\[
\begin{cases}
  x + y - 2z = 6 \\
  2x + 3y - 7z = 16 \\
  5x + 2y + z = 16
\end{cases}
\]

\begin{cases}
x + y - 2z = 6 \\
2x + 3y - 7z = 16\\
5x + 2y + z = 16\\
\end{cases}
\rightarrow
\begin{pmatrix}
1 & 1 & -2 \\
2 & 3 & -7 \\
5 & 2 & 1
\end{pmatrix}
\begin{pmatrix}
    x\\
    y\\
    z
\end{pmatrix}
=
\begin{pmatrix}
6 \\
16\\
16
\end{pmatrix}
\rightarrow
\left[
\begin{array}{ccc|c} 
1 & 1 & -2 & 6 \\ 
2 & 3 & -7 & 16 \\
5 & 2 & 1 & 16
\end{array}
\right]
\rightarrow
\rightarrow
\left[
\begin{array}{ccc|c} 
0 & -1 & 3 & -4 \\ 
0 & 11 & -37 & 48 \\
10 & 4 & 2 & 32
\end{array}
\right]
\rightarrow
\left[
\begin{array}{ccc|c} 
0 & 0 & -4 & 4 \\ 
0 & 11 & -37 & 48 \\
10 & 4 & 2 & 32
\end{array}
\right]
\rightarrow
\left[
\begin{array}{ccc|c} 
0 & 0 & 1 & -1 \\ 
0 & 11 & -37 & 48 \\
5 & 2 & 1 & 16
\end{array}
\right]
\rightarrow
\left[
\begin{array}{ccc|c} 
0 & 0 & 1 & -1 \\ 
0 & 11 & 0 & 11 \\
5 & 2 & 1 & 16
\end{array}
\right]
\rightarrow
\left[
\begin{array}{ccc|c} 
0 & 0 & 1 & -1 \\ 
0 & 1 & 0 & 1 \\
5 & 2 & 1 & 16
\end{array}
\right]
\rightarrow
\left[
\begin{array}{ccc|c} 
0 & 0 & 1 & -1 \\ 
0 & 1 & 0 & 1 \\
5 & 0 & 0 & 15
\end{array}
\right]
\rightarrow
\left[
\begin{array}{ccc|c} 
0 & 0 & 1 & -1 \\ 
0 & 1 & 0 & 1 \\
1 & 0 & 0 & 3
\end{array}
\right]

Thus, $x = 3, y = 1, z = -1$

\item Find the eigenvectors and eigenvalues of a linear operator
\[
A = \begin{pmatrix}
1 & 0 & 3 \\
2 & 1 & 2 \\
3 & 0 & 1
\end{pmatrix}
\]

$det(A-\lambda I) = 0$

$|A-\lambda I| = $
\left|
\begin{array}{ccc} 
1-\lambda & 0 & 3 \\ 
2 & 1-\lambda & 2  \\
3 & 0 & 1-\lambda 
\end{array}
\right|
$= (1-\lambda)^3 + 3(-3(1-\lambda) = (1-\lambda)^3 - 9(1-\lambda) = (1-\lambda)(\lambda^2-2\lambda-8) = (1-\lambda)(\lambda-4)(\lambda+2)=0

Thus, eigenvalues are $\lambda_1 = 1, \lambda_2 = 4, \lambda_3 = -2$


$\lambda_1 = 1 \rightarrow (A - \lambda_1I)
\left[
\begin{array}{c} 
x_1 \\ 
x_2  \\
x_3 
\end{array}
\right] = 0

\left(
\begin{array}{ccc|c} 
0 & 0 & 3 & 0 \\ 
2 & 0 & 2 & 0 \\
3 & 0 & 0 & 0
\end{array}
\right)
\rightarrow
\left(
\begin{array}{ccc|c} 
0 & 0 & 1 & 0 \\ 
1 & 0 & 1 & 0 \\
1 & 0 & 0 & 0
\end{array}
\right)
\rightarrow
\left(
\begin{array}{ccc|c} 
0 & 0 & 1 & 0 \\ 
0 & 0 & 0 & 0 \\
1 & 0 & 0 & 0
\end{array}
\right)

Thus, $x_1 = 0, x_2 \in \mathbb{R}, x_3 = 0$

If $x_2 = 1$, for $\lambda_1 = 1, v_1 = (0,1,0)$

$\lambda_2 = 4 \rightarrow (A - \lambda_2I)
\left[
\begin{array}{c} 
x_1 \\ 
x_2  \\
x_3 
\end{array}
\right] = 0

\left(
\begin{array}{ccc|c} 
-3 & 0 & 3 & 0 \\ 
2 & -3 & 2 & 0 \\
3 & 0 & -3 & 0
\end{array}
\right)
\rightarrow
\left(
\begin{array}{ccc|c} 
-1 & 0 & 1 & 0 \\ 
2 & -3 & 2 & 0 \\
1 & 0 & -1 & 0
\end{array}
\right)
\rightarrow
\left(
\begin{array}{ccc|c} 
-1 & 0 & 1 & 0 \\ 
2 & -3 & 2 & 0 \\
0 & 0 & 0 & 0
\end{array}
\right)
\rightarrow
\left(
\begin{array}{ccc|c} 
-2 & 0 & 2 & 0 \\ 
0 & -3 & 4 & 0 \\
0 & 0 & 0 & 0
\end{array}
\right)
\rightarrow
\left(
\begin{array}{ccc|c} 
-1 & 0 & 1 & 0 \\ 
0 & -3 & 4 & 0 \\
0 & 0 & 0 & 0
\end{array}
\right)
\rightarrow
\begin{cases}
-x_1 + x_3 = 0 \\
-3x_2 + 4x_3 = 0
\end{cases}
\rightarrow
\begin{cases}
x_1 = x_3\\
x_2 = \frac{4}{3}x_3\\
x_3 \in \mathbb{R}
\end{cases}

If $x_3 = 1$, then for $\lambda_2 = 4$, eigenvector $v_2 = (1, \frac{4}{3},1)$

$\lambda_3 = -2 \rightarrow (A - \lambda_3I)
\left[
\begin{array}{c} 
x_1 \\ 
x_2  \\
x_3 
\end{array}
\right] = 0

\left(
\begin{array}{ccc|c} 
3 & 0 & 3 & 0 \\ 
2 & 3 & 2 & 0 \\
3 & 0 & 3 & 0
\end{array}
\right)
\rightarrow
\left(
\begin{array}{ccc|c} 
1 & 0 & 1 & 0 \\ 
0 & 1 & 0 & 0 \\
0 & 0 & 0 & 0
\end{array}
\right)
\rightarrow
\begin{cases}
x_1 + x_3 = 0 \\
x_2 = 0 \\
\end{cases}
\rightarrow
\begin{cases}
x_1 = - x_3 \\
x_2 = 0 \\
x_3 \in \mathbb{R}
\end{cases}

If $x_3 = 1$, then for $\lambda_3 = -2$, eigenvector $v_3 = (-1, 0, 1)$

\item How many integers between 0 and 999 contain

a) exactly one digit 7?

Since there are 3 digits, and only one of them is 7, total number of integers with exactly one seven is $3*9^2 = 243$

b) at least one digit 7?

Since there are 1000 integers between 0 and 999 and $9^3$ integers that does not contain 7, the number of integers that contain and least one 7 is $1000 - 9^3 = 271$


\item How many numbers less than one million can be written using the digits 8 and 9?

Since the largest possible number is 999999, all integers must contain at most 6 digits.

For 1 digit numbers there are $2^1$ possible numbers - 8,9

For 2 digit numbers there are $2^2$ possible numbers - 88,99,89,98

...

For 6 digit numbers there are $2^6$ possible numbers

So the total number of such integers is $2 + 2^2 + 2^3 + 2^4 + 2^5 + 2^6=126$ 


\item How many four-digit numbers greater than 2100 can be formed from the digits
1, 2, 3, 4, 5, 0 if no digit is repeated in the number?

N of such integers = n (numbers starting with 2 and greater than 2100) + n (numbers starting with 3,4, or 5, and greater than 2100) = $(1*4*4*3) + (3*5*4*3) = 48 + 180 = 228$ 


\end{enumerate}
\end{document}


